\section*{Hoofdstuk \ref{ch:g2:mixing-and-dissolving} --- Mengen en oplossen}

\begin{solution}{exo:g2:mixing-and-dissolving:1}
Suiker en zout \omterm{def:g2:mixing-and-dissolving:dissolve}{lossen op}. Zand, krijt en kiezels niet.
\end{solution}

\begin{solution}{exo:g2:mixing-and-dissolving:2}
Een \omterm{def:g2:mixing-and-dissolving:dissolve}{oplossing} is de heldere vloeistof die je krijgt als een vaste stof
in water \omterm{def:g2:mixing-and-dissolving:dissolve}{oplost}. Voorbeelden: suikerwater, zout water (zoete thee is er
ook een).
\end{solution}

\begin{solution}{exo:g2:mixing-and-dissolving:3}
Ja: havervlokken, rozijnen, noten en zaadjes zijn bij elkaar gedaan. Elk
deel is nog te zien en eruit te halen.
\end{solution}

\begin{solution}{exo:g2:mixing-and-dissolving:4}
Nee. Een opgeloste vaste stof laat het water helder, met niets op de
bodem; hier is het water troebel en is er een laagje bezonken.
\end{solution}

\begin{solution}{exo:g2:mixing-and-dissolving:5}
Het is helder (we kunnen erdoorheen kijken) maar niet kleurloos (het is
roze). Het is een \omterm{def:g2:mixing-and-dissolving:dissolve}{oplossing}: een \omterm{def:g2:mixing-and-dissolving:dissolve}{oplossing} mag een kleur hebben.
\end{solution}

\begin{solution}{exo:g2:mixing-and-dissolving:6}
De thee smaakt zoet: de suiker zit er nog in, door het water verspreid in
stukjes die te klein zijn om te zien.
\end{solution}

\begin{solution}{exo:g2:mixing-and-dissolving:7}
De zon laat het water opdrogen; het zout dat in het zeewater opgelost
was, blijft achter.
\end{solution}

\begin{solution}{exo:g2:mixing-and-dissolving:8}
``Opgelost.'' Smelten is wat een ijsblokje vanzelf doet, zonder water;
de suiker heeft het water nodig, en hij zit nog in de thee.
\end{solution}

\begin{solution}{exo:g2:mixing-and-dissolving:9}
Het bekerglas met het zand is troebel en heeft een laagje op de bodem;
het suikerwater is helder, met niets op de bodem.
\end{solution}

\begin{solution}{exo:g2:mixing-and-dissolving:10}
Kraanwater is niet alleen water: er is iets in opgelost. Als het water
opdroogt, blijft die vaste stof achter als het witte kringetje.
\end{solution}

\begin{solution}{exo:g2:mixing-and-dissolving:11}
Ze laat van elk glas een paar druppels opdrogen op een schoon, donker
schaaltje. Het zuivere water laat niets achter; het zoute water laat wit
zout achter.
\end{solution}
